% ===== PRÉAMBULE CANONIQUE — à coller VERBATIM en tête de CHAQUE .tex =====
\documentclass[11pt,a4paper]{article}
\usepackage[utf8]{inputenc}\usepackage[T1]{fontenc}\usepackage{lmodern}
\usepackage[french]{babel}
\usepackage{amsmath,amssymb,mathtools}\usepackage{icomma}
\usepackage[version=4]{mhchem}          % \ce{...} : équations & formules
\usepackage{siunitx}\sisetup{locale=FR} % \qty{}{}, \si{}, \ang{}
\usepackage{chemfig}                    % \chemfig{...} : structures développées
\usepackage{xcolor}\usepackage{colortbl}
\usepackage{tikz}\usepackage{pgfplots}\pgfplotsset{compat=1.18}
\usepackage[most]{tcolorbox}\usepackage{enumitem}\usepackage{array,booktabs}
\usepackage[a4paper,margin=2.2cm]{geometry}
\usetikzlibrary{arrows.meta,calc,angles,quotes,positioning,patterns,backgrounds,shapes.geometric}
\definecolor{primary}{RGB}{222,49,99}\definecolor{primarydark}{RGB}{150,22,55}
\definecolor{defcol}{RGB}{0,114,178}\definecolor{methcol}{RGB}{52,92,148}
\definecolor{excol}{RGB}{0,135,90}\definecolor{attncol}{RGB}{200,40,55}
\tcbset{boxbase/.style={enhanced,breakable,boxrule=0.9pt,arc=2.5pt,
  left=3mm,right=3mm,top=2.3mm,bottom=2.3mm,fonttitle=\bfseries,coltitle=white}}
% --- boîtes TUTORIEL (noms EXACTS, ne pas renommer) ---
\newtcolorbox{cle}[1][]{boxbase,colback=defcol!6,colframe=defcol,
  title={\textsc{À retenir}\ifx\relax#1\relax\relax\else\textmd{ : \itshape #1}\fi}}
\newtcolorbox{methode}[1][]{boxbase,colback=methcol!7,colframe=methcol,sharp corners,
  title={\textsc{Méthode}\ifx\relax#1\relax\relax\else\textmd{ --- \itshape #1}\fi}}
\newtcolorbox{exemplebox}[1][]{boxbase,colback=excol!5,colframe=excol,
  title={\textsc{Exemple}\ifx\relax#1\relax\relax\else\textmd{ : \itshape #1}\fi}}
\newtcolorbox{piege}[1][]{boxbase,colback=attncol!6,colframe=attncol,
  title={\textsc{Piège}\ifx\relax#1\relax\relax\else\textmd{ : \itshape #1}\fi}}
\newtcolorbox{remarque}[1][]{boxbase,colback=black!4,colframe=black!45,
  title={\textsc{Remarque}\ifx\relax#1\relax\relax\else\textmd{ : \itshape #1}\fi}}
% --- boîtes EXERCICE / CORRIGÉ (noms EXACTS, auto-comptées) ---
\newtcolorbox[auto counter]{exo}[1][]{boxbase,colback=primary!4,colframe=primary,
  title={\textsc{Exercice}~\thetcbcounter\ifx\relax#1\relax\relax\else\textmd{ --- \itshape #1}\fi}}
\newtcolorbox[auto counter]{corrige}[1][]{boxbase,colback=excol!4,colframe=excol,
  title={\textsc{Corrigé}~\thetcbcounter\ifx\relax#1\relax\relax\else\textmd{ --- \itshape #1}\fi}}
\newcommand{\R}{\mathbb{R}}\newcommand{\N}{\mathbb{N}}\newcommand{\Z}{\mathbb{Z}}\newcommand{\ds}{\displaystyle}
% (un \usepackage{fancyhdr}+en-tête est possible ; il est ignoré côté web)
% =========================================================================
\begin{document}
\sisetup{per-mode=symbol}
\begin{center}\small\itshape On rappelle : $\log 2 \approx 0{,}30$,\quad $\log 4 = 2\log 2 \approx 0{,}60$,\quad $\log 2{,}5 \approx 0{,}40$.\end{center}

\begin{exo}[mélange de trois apports d'acides forts]
On prépare une solution en mélangeant : \qty{50}{mL} contenant à la fois \ce{HCl} à \qty{0,10}{mol\per L} et
\ce{HBr} à \qty{0,10}{mol\per L}, puis \qty{50}{mL} de \ce{HCl} à \qty{0,60}{mol\per L}. (\ce{HCl} et \ce{HBr}
sont des acides forts ; volume total \qty{100}{mL}.)
\begin{enumerate}[label=\alph*)]
  \item Calcule le nombre de moles de \ce{H3O+} apportées par chaque source.
  \item En déduis $[\ce{H3O+}]$, puis le pH de la solution.
\end{enumerate}
\end{exo}

\begin{exo}[mélange de trois bases fortes]
On mélange \qty{25}{mL} de \ce{NaOH} à \qty{1,0}{mol\per L}, \qty{25}{mL} de \ce{Ba(OH)2} à
\qty{0,10}{mol\per L} et \qty{50}{mL} de \ce{KOH} à \qty{0,20}{mol\per L} (volume total \qty{100}{mL}).
\begin{enumerate}[label=\alph*)]
  \item Calcule le nombre de moles d'ions \ce{OH-} apportées par chaque base (attention à \ce{Ba(OH)2}).
  \item En déduis $[\ce{OH-}]$, le pOH, puis le pH.
\end{enumerate}
\end{exo}

\begin{exo}[diluer pour tripler le degré d'ionisation]
On dispose de \qty{1,0}{L} d'acide acétique \ce{CH3COOH} ($\mathrm{p}K_a = 4{,}75$) à \qty{0,10}{mol\per L}. On
veut \textbf{tripler} son degré d'ionisation $\alpha$ par dilution.
\begin{enumerate}[label=\alph*)]
  \item D'après la loi d'Ostwald, par quel facteur faut-il diviser la concentration ?
  \item Quel volume final faut-il atteindre, et donc quel volume d'eau faut-il ajouter ?
\end{enumerate}
\end{exo}

\begin{exo}[préparer une solution à partir d'un acide fort concentré]
On prélève \qty{40}{mL} d'une solution d'\ce{HCl} à \qty{5,0}{mol\per L} que l'on verse dans une fiole jaugée,
puis on complète à \qty{1,0}{L} avec de l'eau distillée.
\begin{enumerate}[label=\alph*)]
  \item Quelle quantité de matière d'\ce{HCl} contient la solution finale ?
  \item Calcule sa concentration, puis son pH.
\end{enumerate}
\end{exo}

\begin{exo}[capstone : même concentration, acidités différentes]
On compare deux solutions préparées à la \textbf{même} concentration $C = \qty{0,010}{mol\per L}$ : l'une
d'acide fort \ce{HCl}, l'autre d'acide faible \ce{HA} ($\mathrm{p}K_a = 4{,}8$).
\begin{enumerate}[label=\alph*)]
  \item Calcule le pH de la solution de \ce{HCl}.
  \item Calcule le pH de la solution de \ce{HA}.
  \item Laquelle est la plus acide, et par quel facteur les concentrations en \ce{H3O+} diffèrent-elles ?
\end{enumerate}
\end{exo}

\end{document}
